\chapter{Comprendre les fonctions}
\label{corriges:08}
\section*{Exercice 8.1 -- Images et antécédents}
$f(-3)=5$, $f(0)=-4$, $f(2)=0$.
$f(x)=5$ donne $x^2=9$, donc $x=-3$ ou $3$.
$f(x)=-4$ donne $x=0$.
$f(x)=-5$ donne $x^2=-1$, impossible dans les réels.

\section*{Exercice 8.2 -- Lire un tableau avec prudence}
L'image de $-1$ est $0$. Les nombres $-2$ et $2$ sont des antécédents de $3$. La valeur de $g(1{,}5)$ n'est pas déterminée par le tableau. Même si $x^2-1$ est une formule compatible, rien n'affirme que c'est la fonction étudiée. On ne peut pas interpoler sans hypothèse.

\section*{Exercice 8.3 -- Du programme à la formule}
$h(x)=-3x+7$. L'antécédent de $19$ vérifie $-3x+7=19$, donc $-3x=12$ et $x=-4$. Vérification : $h(-4)=12+7=19$.

\section*{Exercice 8.4 -- Modéliser une aire}
La somme des deux côtés vaut $15$, donc $A(x)=x(15-x)$ pour $0<x<15$.
$A(4)=4\times11=44\unite{cm^2}$. L'autre longueur $x=11$ donne $11\times4=44$ : c'est le même rectangle avec les rôles des côtés échangés.
